\documentclass[10pt,a4paper]{article}
\usepackage[margin=2.2cm]{geometry}
\usepackage{fontspec}
\setmonofont{DejaVu Sans Mono}[Scale=0.84]
\newfontfamily\unifont{DejaVu Serif}[Scale=0.9]
\usepackage{amsmath,amssymb,mathtools}
\setcounter{MaxMatrixCols}{20}
\usepackage{longtable,booktabs,array}
\usepackage{fancyvrb}
\usepackage[dvipsnames]{xcolor}
\usepackage[colorlinks=true,linkcolor=NavyBlue,urlcolor=NavyBlue]{hyperref}
\usepackage{microtype}
\usepackage{needspace}
\usepackage{tocloft}\setlength{\cftbeforesecskip}{0.45em}
\setlength{\parindent}{0pt}\setlength{\parskip}{0.55em}
\setlength{\emergencystretch}{8em}
\makeatletter
\renewcommand\section{\@startsection{section}{1}{\z@}{-3.5ex \@plus -1ex \@minus -.2ex}{2.3ex \@plus.2ex}{\normalfont\Large\bfseries\raggedright}}
\renewcommand\subsection{\@startsection{subsection}{2}{\z@}{-3.25ex\@plus -1ex \@minus -.2ex}{1.5ex \@plus .2ex}{\normalfont\large\bfseries\raggedright}}
\renewcommand\subsubsection{\@startsection{subsubsection}{3}{\z@}{-3.25ex\@plus -1ex \@minus -.2ex}{1.5ex \@plus .2ex}{\normalfont\normalsize\bfseries\raggedright}}
\makeatother
\renewcommand{\arraystretch}{1.15}
\allowdisplaybreaks
\begin{document}
\title{The radion of 7+1 spacetime: symbols and their definitions\\[0.4em]\large A dictionary of every symbol used in the document on the radion, with its defining equation and its name in the notebook 7+1.nb}
\author{Repository: github.com/once-ere/7plus1 — prepared for Nash\\ Notebook: 7+1.nb (Wolfram Language 15.0.1)\\ Document generated by tools/make\_docs.py on 2026-10-28}
\date{}
\maketitle
\begin{abstract}
For each symbol this dictionary gives the mathematical notation, the name used in the Mathematica notebook 7+1.nb (where there is one; the names of section 26 begin with \texttt{rd}; "—" means that the notebook has no separate name), the defining equation and its meaning, grouped by topic in the order of the document. Longer defining equations follow each table as numbered displayed formulas. Section 11 lists the symbols that look alike, section 12 every Mathematica name of section 26. Every equation needed to define a symbol is written in this document.
\end{abstract}
\tableofcontents
\Needspace{12\baselineskip}\section*{1. Conventions}\addcontentsline{toc}{section}{1. Conventions}
\begin{itemize}
\item Units: $\hbar=c=1$; masses and inverse lengths in GeV; $\hbar c=1.973269804\times10^{-16}$ GeV m (Mathematica \texttt{tbHbarCm}).
\item Coordinates $X=(x_0,x_1,\dots,x_7)$ (Mathematica \texttt{X} $=$ \texttt{\{x0, …, x7\}}): $x_0$ is time (the direction of the octonion unit $e_0=1$), $x_1,x_2,x_3$ ordinary space, $x_4,\dots,x_7$ the four extra space-like coordinates: on the 4-torus periodic with period $L_c$, on the 4-sphere the angles $x_4,x_5,x_6\in[0,\pi]$, $x_7\in[0,2\pi)$.
\item Indices: $A,B=0,\dots,7$; $\mu,\nu=0,\dots,3$ (ordinary spacetime, the directions of the world brane); $i,j=4,\dots,7$ (the extra space); $a,b=1,\dots,4$ (frame indices of the extra space, or the coordinates $X_a$ of $\mathbb R^5$ when $a,b\le5$). A prime is $d/dx_0$ (cosmology) or $d/dy$ (warped metrics); $\ln$ is the natural logarithm.
\item Curvature (mostly-minus signature): $R^\rho{}_{\sigma\mu\nu}=\partial_\mu\Gamma^\rho{}_{\nu\sigma}-\partial_\nu\Gamma^\rho{}_{\mu\sigma}+\Gamma^\rho{}_{\mu\lambda}\Gamma^\lambda{}_{\nu\sigma}-\Gamma^\rho{}_{\nu\lambda}\Gamma^\lambda{}_{\mu\sigma}$, $R_{\sigma\nu}=R^\rho{}_{\sigma\rho\nu}$, $R=g^{\sigma\nu}R_{\sigma\nu}$, $G^A{}_B=R^A{}_B-\frac12\delta^A_BR$; field equations $G^A{}_B-\Lambda\delta^A_B=\kappa T^A{}_B$.
\end{itemize}
\textbf{The signature: $\{7,1,0\}$ and $\mathrm{diag}(+1,-1,\dots,-1)$.} The 7+1 spacetime has one time-like dimension ($x_0$, the direction of the octonion unit $e_0=1$) and seven space-like dimensions ($x_1,x_2,x_3$ ordinary space, $x_4,\dots,x_7$ auxiliary space). Two different objects describe this, and neither may be replaced by the other. (i) The tensor package xAct, used in section 12 of 7+1.nb exactly as in the author’s notebook Einstein-Lovelock-7+1-Nash.nb, declares the metric with \texttt{DefMetric[\{7, 1, 0\}, g71[-a, -b], CD, …]}: the list $\{7,1,0\}$ declares seven space-like dimensions, one time-like dimension and no degenerate dimension — the numbers of entries $+1$ and $-1$ of $\eta_{71}=\mathrm{diag}(-1,+1,\dots,+1)$ (\texttt{η71}), the same spacetime written in the sign convention of 4+4.nb (space-like $=+1$, time-like $=-1$). (ii) The matrix $\eta=-\eta_{71}=\mathrm{diag}(+1,-1,\dots,-1)$ (\texttt{η17}) is simply the squares of the octonion basis: $e_0^2=+1$ for the unit $e_0=1$ (time) and $e_i^2=-1$ for the seven imaginary units $e_1,\dots,e_7$ (space); more generally it is the real part of the product of two basis units, and it gives the quadratic form of the octonions:
\begin{gather*}
\eta_{AB}=\mathrm{Re}(e_Ae_B),\qquad e_0^2=+1,\qquad e_i^2=-1\ \ (i=1,\dots,7),\\ \mathrm{Re}(x\,x)=x_0^2-x_1^2-\cdots-x_7^2=\eta_{AB}\,x^Ax^B\qquad (x=x^Ae_A\in\mathbb O).
\end{gather*}
$\eta$ gives the signs of the Clifford relation $\Gamma_A\Gamma_B+\Gamma_B\Gamma_A=2\,\eta_{AB}\,1_{16}$ of the real 16×16 Dirac matrices: time squares to $+1_{16}$ and each of the seven space directions to $-1_{16}$, exactly like the octonion units (in the notation $Cl(p,q)$, with $p$ generators squaring to $+1$ and $q$ to $-1$, the algebra $Cl(1,7)\cong\mathbb R(16)$). It also gives the signs of the metric components: $g_{00}>0$ along the time-like direction, negative squares along the space-like directions. Every component computed by 7+1.nb follows from the component matrix of the metric; the declared list fixes for xAct the sign of the determinant of the metric, \texttt{SignDetOfMetric[g71]} $=-1$, which is the same in both sign conventions, since $\det \eta=\det \eta_{71}=-1$. 7+1.nb checks \texttt{SignatureOfMetric[g71]} $=\{7,1,0\}$, \texttt{SignDetOfMetric[g71]} $=-1$, $\mathrm{Re}(e_Ae_B)=\eta_{AB}$ and $\mathrm{Re}(x\,x)=\eta_{AB}x^Ax^B$.

\Needspace{12\baselineskip}\section*{2. Coordinates, metrics and the two extra spaces}\addcontentsline{toc}{section}{2. Coordinates, metrics and the two extra spaces}
{\footnotesize
\begin{longtable}{>{\raggedright\arraybackslash}p{\dimexpr 0.1700\linewidth-2\tabcolsep\relax}>{\raggedright\arraybackslash}p{\dimexpr 0.2200\linewidth-2\tabcolsep\relax}>{\raggedright\arraybackslash}p{\dimexpr 0.3100\linewidth-2\tabcolsep\relax}>{\raggedright\arraybackslash}p{\dimexpr 0.3000\linewidth-2\tabcolsep\relax}}
\toprule
\textbf{symbol} & \textbf{Mathematica} & \textbf{definition} & \textbf{meaning} \\
\midrule\endfirsthead
\toprule
\textbf{symbol} & \textbf{Mathematica} & \textbf{definition} & \textbf{meaning} \\
\midrule\endhead
\endfoot
\bottomrule\endlastfoot
$x$, $y$ & \texttt{\{x0, …, x3\}}, \texttt{\{x4, …, x7\}} & $x=(x_0,\dots,x_3)$, $y=(x_4,\dots,x_7)$ & a point of ordinary spacetime; a point of the extra space \\
$G_{AB}$, $ds^2$ & — & (2.1) & the eight-dimensional metric of the radion ansatz \\
$g_{E\mu\nu}$ & \texttt{\{rdLp\textasciicircum{}2, -rdSa\textasciicircum{}2, …\}} (homogeneous) & (2.1) & the four-dimensional Einstein-frame metric \\
$N$, $a$ & \texttt{rdLp[x0]}, \texttt{rdSa[x0]} & $g_E=\mathrm{diag}(N^2,-a^2,-a^2,-a^2)$ & lapse and scale factor of the homogeneous checks \\
$\gamma_{ij}$ & \texttt{rdB0\textasciicircum{}2 \{1,1,1,1\}}, \texttt{rdR0\textasciicircum{}2 rdGam4} & $b_0^2\delta_{ij}$ ($T^4$); $r_0^2\hat\gamma_{ij}$ ($S^4$) & the fixed reference metric of the extra space \\
$\hat\gamma$ & \texttt{rdGam4} & (2.2) & the metric of the unit 4-sphere \\
$T^4$ & — & $\mathbb R^4/(L_c\mathbb Z)^4$ & the 4-torus (the auxiliary space of 7+1.nb) \\
$S^4$ & — & (2.3) & the 4-sphere (for comparison) \\
$L_c$ & \texttt{rdLc} & coordinate period & period of each torus coordinate \\
$b$, $b_0$ & \texttt{b4[x0]} (section 11), \texttt{rdB0} & $b=b_0e^\beta$ & size of the torus; its reference value \\
$r$, $r_0$ & \texttt{rdR0} ($r_0$) & $r=r_0e^\beta$ & radius of the sphere; its reference value \\
$V$, $V_0$ & \texttt{v0} (argument of \texttt{rdMini}) & (2.4) & proper volume of the extra space; at $\beta=0$ \\
$V_\perp$, $\ell_\perp$ & \texttt{rIIVp}, \texttt{tbLp} & $V_\perp=M_{\rm Pl}^2/M_8^6$, $\ell_\perp=V_\perp^{1/4}$ & volume of the torus today; its proper period \\
$r_{S^4}$ & \texttt{rdRS4} & $(3V_\perp/(8\pi^2))^{1/4}$ & radius of a sphere with the volume $V_\perp$ \\
$X_a$ & \texttt{rdEmb} & (2.3) & embedding of $S^4$ in $\mathbb R^5$, $a=1,\dots,5$ \\
$R_{\rm int}$ & \texttt{rint} (argument of \texttt{rdMini}) & $0$ ($T^4$), $12/r^2$ ($S^4$) & scalar curvature of the extra space itself (positive for the sphere) \\
$R$, $R_E$ & \texttt{rdScal[g, X]} & section 1 & scalar curvature of $G$; of $g_E$ \\
$R_{ij}$ & \texttt{rdRic}, \texttt{rdRicS4} & section 1 & Ricci tensor \\*
\end{longtable}}
\begin{align*}
ds^2&=G_{AB}dX^AdX^B=e^{-4\beta(x)}g_{E\mu\nu}(x)dx^\mu dx^\nu-e^{2\beta(x)}\gamma_{ij}(y)dy^idy^j,&&(2.1) \\
\hat\gamma&=\mathrm{diag}\big(1,\ \sin^2x_4,\ \sin^2x_4\sin^2x_5,\ \sin^2x_4\sin^2x_5\sin^2x_6\big),&&(2.2) \\
X_a&=r\,\big(\sin x_4\sin x_5\sin x_6\cos x_7,\ \sin x_4\sin x_5\sin x_6\sin x_7,\ \sin x_4\sin x_5\cos x_6,\ \sin x_4\cos x_5,\ \cos x_4\big),&&(2.3) \\
V&=\int d^4y\sqrt{\det(e^{2\beta}\gamma)}=V_0e^{4\beta},\qquad V_0=b_0^4L_c^4\ (T^4),\quad V_0=\frac{8\pi^2}3r_0^4\ (S^4).&&(2.4)
\end{align*}
\Needspace{12\baselineskip}\section*{3. The radion, the Einstein frame and the reduced action}\addcontentsline{toc}{section}{3. The radion, the Einstein frame and the reduced action}
{\footnotesize
\begin{longtable}{>{\raggedright\arraybackslash}p{\dimexpr 0.1700\linewidth-2\tabcolsep\relax}>{\raggedright\arraybackslash}p{\dimexpr 0.2200\linewidth-2\tabcolsep\relax}>{\raggedright\arraybackslash}p{\dimexpr 0.3100\linewidth-2\tabcolsep\relax}>{\raggedright\arraybackslash}p{\dimexpr 0.3000\linewidth-2\tabcolsep\relax}}
\toprule
\textbf{symbol} & \textbf{Mathematica} & \textbf{definition} & \textbf{meaning} \\
\midrule\endfirsthead
\toprule
\textbf{symbol} & \textbf{Mathematica} & \textbf{definition} & \textbf{meaning} \\
\midrule\endhead
\endfoot
\bottomrule\endlastfoot
$\beta$ & \texttt{rdBt[x0]} (field), \texttt{rdBe} (number) & $\ln(b/b_0)$ or $\ln(r/r_0)$ & the radion (dimensionless) \\
$\phi$ & — & (3.1) & the canonical radion (mass dimension 1) \\
$\alpha$, $\beta$ (the note) & — & $ds^2=e^{2\alpha\phi}g+e^{2\beta\phi}\gamma$ & the constants of the note’s ansatz; $\beta\phi$ of the note is our $\beta(x)$ \\
$\beta_n$ & — & (3.3) & the constant $\beta$ for $n$ extra dimensions \\
$\varphi$ & \texttt{rdPh[x0]} & $\beta=\beta_n\varphi$ & the radion in the normalization of the note (dimensionless) \\
$n$ & \texttt{n} in \texttt{rdKinOK}, \texttt{rdKinC} & $D=4+n$ & number of extra dimensions \\
$c_n$ & \texttt{rdKinC[n]}, \texttt{rdKinCs} & (3.2) & coefficient of the radion kinetic term \\
$\kappa$ & \texttt{rdKap} & $8\pi G_8=M_8^{-6}$ & eight-dimensional gravitational coupling \\
$M_8$ & \texttt{rIIM8} & 10 TeV & eight-dimensional Planck scale of the examples \\
$M_{\rm Pl}$ & \texttt{rIIMPl}, \texttt{rdMP2} ($M_{\rm Pl}^2$) & (3.1) & reduced four-dimensional Planck mass, $2.435\times10^{18}$ GeV \\
$G_4$, $G_8$ & — & $8\pi G_4=M_{\rm Pl}^{-2}$, $8\pi G_8=\kappa$ & Newton constants \\
$\Lambda$ & \texttt{rdLam} & section 1 & bulk cosmological constant ($\Lambda>0$: de Sitter) \\
$U$, $U_0$, $U_1$, $K_U$ & \texttt{rdU}, \texttt{rdU0}, \texttt{rdU1}, \texttt{rdKU}, \texttt{rdUq} & (4.1) & stabilizing energy density of the bulk and its expansion \\
$S$ & — & (3.4) & the action \\
$\Omega$ & \texttt{rdOmg}, \texttt{rdOmE} & $h=\Omega^2\eta$ & conformal factor (section 6) \\*
\end{longtable}}
\begin{align*}
\phi&=\sqrt{12}\,M_{\rm Pl}\,\beta,\qquad M_{\rm Pl}^2=\frac{V_0}\kappa,\qquad 6M_{\rm Pl}^2(\partial\beta)^2=\frac12(\partial\phi)^2,&&(3.1) \\
-\frac1{2\kappa}\int d^ny\sqrt{|G|}R&=\sqrt{|g_E|}\Big[-\frac{M_{\rm Pl}^2}2R_E+\frac{M_{\rm Pl}^2}2c_n(\partial\beta)^2\Big]+\text{(total derivative)},&&(3.2) \\
c_n&=\frac{n(n+2)}2, \\
2\alpha+n\beta_n&=0;\qquad\text{the note: }\beta_n=\sqrt{\frac2{n(n+2)}}\ \ (c_n\beta_n^2=1);&&(3.3) \\
&\qquad\text{canonical: }\beta_n=\frac1{\sqrt{n(n+2)}}\ \ (c_n\beta_n^2=\tfrac12), \\
S&=-\frac1{2\kappa}\int d^8X\sqrt{|G|}(R+2\Lambda)-\int d^8X\sqrt{|G|}\,U+S_{\rm brane}+S_{\rm matter}.&&(3.4)
\end{align*}
\Needspace{12\baselineskip}\section*{4. The potential and the mass}\addcontentsline{toc}{section}{4. The potential and the mass}
{\footnotesize
\begin{longtable}{>{\raggedright\arraybackslash}p{\dimexpr 0.1700\linewidth-2\tabcolsep\relax}>{\raggedright\arraybackslash}p{\dimexpr 0.2200\linewidth-2\tabcolsep\relax}>{\raggedright\arraybackslash}p{\dimexpr 0.3100\linewidth-2\tabcolsep\relax}>{\raggedright\arraybackslash}p{\dimexpr 0.3000\linewidth-2\tabcolsep\relax}}
\toprule
\textbf{symbol} & \textbf{Mathematica} & \textbf{definition} & \textbf{meaning} \\
\midrule\endfirsthead
\toprule
\textbf{symbol} & \textbf{Mathematica} & \textbf{definition} & \textbf{meaning} \\
\midrule\endhead
\endfoot
\bottomrule\endlastfoot
$V(\beta)$ & \texttt{rdV[curv, flux]} & (4.2) & Einstein-frame potential of the radion \\
$\sigma$ & \texttt{curv} (argument of \texttt{rdV}) & $0$ ($T^4$), $1$ ($S^4$) & switch of the curvature term \\
$F$, $\rho_F$, $q$ & \texttt{rdFx}, \texttt{rdRhoF}, \texttt{rdQ2} ($q^2$) & (4.3) & flux term; flux energy density; flux quantum (illustration) \\
$W(a,b)$ & argument of \texttt{feStress} & (4.3) & comoving energy of a source \\
$\varepsilon$, $p_a$, $p_b$ & \texttt{rdEps}, \texttt{rdPa}, \texttt{rdPb} & (4.3) & energy density, pressures along ordinary and extra space \\
$H_a$, $H_b$ & \texttt{Ha}, \texttt{Hb} (section 20) & $a'/a$, $b'/b$ & expansion rates; $H_b=\beta'$ \\
$m_r$, $m_r^2$ & \texttt{rdM2L} (\texorpdfstring{\ensuremath{\Lambda}}{Λ}-only case) & (4.4) & mass of the radion (10 GeV in the examples of 7+1.nb) \\
$H$ & \texttt{rdHE2} ($H^2$) & $H^2=V(0)/(3M_{\rm Pl}^2)$ & Hubble rate of four-dimensional de Sitter space \\
$\Phi$, $m_8$ & \texttt{rdPhi[x4, x5]}, \texttt{rdMm} & (4.5) & a bulk scalar and its mass (the note’s stabilizer) \\*
\end{longtable}}
\begin{align*}
U(\beta)&=U_0+U_1\beta+\frac{K_U}2\beta^2,&&(4.1) \\
V(\beta)&=M_{\rm Pl}^2\Big[e^{-4\beta}\big(\Lambda+\kappa U(\beta)\big)-\frac{6\sigma}{r_0^2}e^{-6\beta}+F\,e^{-12\beta}\Big],&&(4.2) \\
\varepsilon&=\frac W{a^3b^4},\quad p_a=-\frac{\partial_aW}{3a^2b^4},\quad p_b=-\frac{\partial_bW}{4a^3b^3};\qquad W=\frac{a^3b^4q^2}{2b^8}:\ \rho_F=\frac{q^2}{2b^8},\ F=\kappa\rho_F(b_0),&&(4.3) \\
m_r^2&=\frac{d^2V}{d\phi^2}\Big|_0=\frac{V''(0)}{12M_{\rm Pl}^2},&&(4.4) \\
&\sqrt{|G|}\Big(\frac12G^{AB}\partial_A\Phi\partial_B\Phi-\frac12m_8^2\Phi^2\Big)=-\sqrt{|g_E|}\sqrt\gamma\,e^{-4\beta}\Big(\frac12e^{-2\beta}|\nabla\Phi|_\gamma^2+\frac12m_8^2\Phi^2\Big).&&(4.5)
\end{align*}
\Needspace{12\baselineskip}\section*{5. The radion and the world brane}\addcontentsline{toc}{section}{5. The radion and the world brane}
{\footnotesize
\begin{longtable}{>{\raggedright\arraybackslash}p{\dimexpr 0.1700\linewidth-2\tabcolsep\relax}>{\raggedright\arraybackslash}p{\dimexpr 0.2200\linewidth-2\tabcolsep\relax}>{\raggedright\arraybackslash}p{\dimexpr 0.3100\linewidth-2\tabcolsep\relax}>{\raggedright\arraybackslash}p{\dimexpr 0.3000\linewidth-2\tabcolsep\relax}}
\toprule
\textbf{symbol} & \textbf{Mathematica} & \textbf{definition} & \textbf{meaning} \\
\midrule\endfirsthead
\toprule
\textbf{symbol} & \textbf{Mathematica} & \textbf{definition} & \textbf{meaning} \\
\midrule\endhead
\endfoot
\bottomrule\endlastfoot
$W$ (brane) & — & $y=0$ ($T^4$), $x_4=0$ ($S^4$) & the world brane \\
$h_{\mu\nu}$ & — & (5.1) & induced metric on the brane \\
$T^\mu{}_\mu$ & — & trace of $T^{\mu\nu}$ & trace of the energy–momentum tensor of brane matter \\
$M$, $M'$ & \texttt{rdMs} & masses & point masses on the brane \\
$g$ & — & (5.2) & coupling of the radion to $T^\mu{}_\mu$ \\
$\alpha$ (fifth force) & — & (5.2) & strength of the radion force relative to gravity \\
$\lambda$ (Yukawa) & — & $1/m_r$ & range of a Yukawa force \\
$r$ (distance) & — & $|\vec x-\vec x\,'|$ & distance between two masses \\*
\end{longtable}}
\begin{align*}
h_{\mu\nu}&=e^{-4\beta}g_{E\mu\nu}\quad(n=4;\ \text{in general }e^{-n\beta}g_{E\mu\nu}),&&(5.1) \\
\mathcal L_{\rm int}&=g\,\phi\,T^\mu{}_\mu,\quad g=\frac1{\sqrt3M_{\rm Pl}}\ \Big(\text{in general }\frac1{M_{\rm Pl}}\sqrt{\frac n{2(n+2)}}\Big),\quad \alpha=2(gM_{\rm Pl})^2=\frac n{n+2}.&&(5.2)
\end{align*}
\Needspace{12\baselineskip}\section*{6. The Higgs boson and the radion}\addcontentsline{toc}{section}{6. The Higgs boson and the radion}
{\footnotesize
\begin{longtable}{>{\raggedright\arraybackslash}p{\dimexpr 0.1700\linewidth-2\tabcolsep\relax}>{\raggedright\arraybackslash}p{\dimexpr 0.2200\linewidth-2\tabcolsep\relax}>{\raggedright\arraybackslash}p{\dimexpr 0.3100\linewidth-2\tabcolsep\relax}>{\raggedright\arraybackslash}p{\dimexpr 0.3000\linewidth-2\tabcolsep\relax}}
\toprule
\textbf{symbol} & \textbf{Mathematica} & \textbf{definition} & \textbf{meaning} \\
\midrule\endfirsthead
\toprule
\textbf{symbol} & \textbf{Mathematica} & \textbf{definition} & \textbf{meaning} \\
\midrule\endhead
\endfoot
\bottomrule\endlastfoot
$H$, $\varphi_H$ & \texttt{rdPhE} ($\varphi_H$) & $H^\dagger H=\varphi_H^2/2$, $\varphi_H=v+h$ & Higgs doublet; its radial field \\
$v$ & \texttt{vEW}, \texttt{rdVv} & $(\sqrt2G_F)^{-1/2}=246.22$ GeV & Higgs vacuum value \\
$h$, $r$ & \texttt{rdHh[x0, x1]}, \texttt{rdRr[x0, x1]} & section 6 & bare Higgs and radion fields (canonical $r$) \\
$\lambda$ (Higgs) & \texttt{rdLamH} & $V=\frac\lambda4(\varphi_H^2-v^2)^2$ & Higgs quartic coupling \\
$\xi$ & \texttt{rdXi} & (6.1) & curvature coupling of the Higgs ($\frac16$: conformal) \\
$\Lambda_r$, $\gamma$ & \texttt{rdO1} ($1/\Lambda_r$), \texttt{rdGam} & $\Omega=1+r/\Lambda_r+\dots$, $\gamma=v/\Lambda_r$ & coupling scale of the radion; mixing parameter \\
$m_{h0}$, $m_{r0}$ & \texttt{rdMh0}, \texttt{rdMr0} & $m_{h0}^2=2\lambda v^2$ & bare Higgs and radion masses \\
$c$ & \texttt{rdc4} & (6.2) & coefficient of $\Omega\Box\Omega$ \\
$K$, $K_{hr}$, $K_{rr}$ & \texttt{rdKm}, \texttt{rdKhr}, \texttt{rdKrr} & (6.3) & kinetic matrix of $(h,r)$ \\
$S$ (shift), $Z$ & \texttt{rdSm}, \texttt{rdZ}, \texttt{rdZ2} ($Z^2$) & (6.4) & kinetic shift; normalization \\
$M'$ & \texttt{rdMp} & (6.5) & mass matrix of $(h',r')$ \\
$Q(\theta)$, $\theta$ & \texttt{rdQ[th]}, \texttt{rdth} & (6.5) & rotation; mixing angle \\
$h_m$, $r_m$ & — & (6.5) & mass eigenstates \\
$C_h$, $C_r$ & — & (6.5) & coefficients of $h_m$ in $h$ and $r$ \\
$b_{\rm EM}$, $F_{\rm SM}$, $b_{\rm QCD}$, $F_{\rm heavy}$ & — & $-\frac{11}3$, $-6.5$, $7$, $-\frac43$ (the note) & trace-anomaly and loop coefficients \\
$\Gamma$ & \texttt{rdHRnum}, \texttt{rdHR71} & (6.6) & decay rates of the Higgs boson \\*
\end{longtable}}
\begin{align*}
\mathcal L&=\sqrt{|h|}\Big[\frac12h^{\mu\nu}\partial_\mu\varphi_H\partial_\nu\varphi_H-\frac\lambda4(\varphi_H^2-v^2)^2+\frac\xi2R(h)\varphi_H^2\Big]+\frac12(\partial r)^2-\frac12m_{r0}^2r^2,&&(6.1) \\
h&=\Omega^2\eta:\quad \sqrt{|h|}R(h)=c\,\Omega\Box\Omega,\quad c=-6,\quad \Box=\partial_0^2-\nabla^2,&&(6.2) \\
\mathcal L^{(2)}&=\frac12\begin{pmatrix}\partial h&\partial r\end{pmatrix}K\begin{pmatrix}\partial h\\ \partial r\end{pmatrix}-\frac12m_{h0}^2h^2-\frac12m_{r0}^2r^2,\quad K=\begin{pmatrix}1&6\xi\gamma\\ 6\xi\gamma&1+6\xi\gamma^2\end{pmatrix},&&(6.3) \\
\begin{pmatrix}h\\ r\end{pmatrix}&=S\begin{pmatrix}h'\\ r'\end{pmatrix},\quad S=\begin{pmatrix}1&-6\xi\gamma/Z\\ 0&1/Z\end{pmatrix},\quad S^{\mathsf T}KS=1\iff Z^2=1+6\xi\gamma^2(1-6\xi),&&(6.4) \\
M'&=S^{\mathsf T}\mathrm{diag}(m_{h0}^2,m_{r0}^2)S,\quad \begin{pmatrix}h'\\ r'\end{pmatrix}=Q(\theta)\begin{pmatrix}h_m\\ r_m\end{pmatrix},\quad Q(\theta)=\begin{pmatrix}\cos\theta&-\sin\theta\\ \sin\theta&\cos\theta\end{pmatrix},&&(6.5) \\
&\tan2\theta=\frac{12\gamma\xi Zm_{h0}^2}{m_{r0}^2-m_{h0}^2(Z^2-36\xi^2\gamma^2)},\quad C_h=\cos\theta-\frac{6\xi\gamma}Z\sin\theta,\quad C_r=\frac{\sin\theta}Z, \\
\frac{\Gamma(h\to\gamma\gamma)}{\Gamma_{\rm SM}}&=\Big|C_h+C_r\gamma\frac{b_{\rm EM}}{F_{\rm SM}}\Big|^2,\qquad \frac{\Gamma(h\to gg)}{\Gamma_{\rm SM}}=\Big|C_h+C_r\gamma\frac{b_{\rm QCD}}{F_{\rm heavy}}\Big|^2.&&(6.6)
\end{align*}
\Needspace{12\baselineskip}\section*{7. Kaluza–Klein levels on \texorpdfstring{$T^4$}{T4} and \texorpdfstring{$S^4$}{S4}}\addcontentsline{toc}{section}{7. Kaluza–Klein levels on \texorpdfstring{$T^4$}{T4} and \texorpdfstring{$S^4$}{S4}}
{\footnotesize
\begin{longtable}{>{\raggedright\arraybackslash}p{\dimexpr 0.1700\linewidth-2\tabcolsep\relax}>{\raggedright\arraybackslash}p{\dimexpr 0.2200\linewidth-2\tabcolsep\relax}>{\raggedright\arraybackslash}p{\dimexpr 0.3100\linewidth-2\tabcolsep\relax}>{\raggedright\arraybackslash}p{\dimexpr 0.3000\linewidth-2\tabcolsep\relax}}
\toprule
\textbf{symbol} & \textbf{Mathematica} & \textbf{definition} & \textbf{meaning} \\
\midrule\endfirsthead
\toprule
\textbf{symbol} & \textbf{Mathematica} & \textbf{definition} & \textbf{meaning} \\
\midrule\endhead
\endfoot
\bottomrule\endlastfoot
$\Psi$, $\psi_n$, $Y_n$ & — & $\Psi=\sum_n\psi_n(x)Y_n(y)$ & a bulk field, its four-dimensional modes and the eigenfunctions \\
$\Delta_\gamma$ & — & Laplacian of $\gamma$ & Laplacian of the extra space \\
$\mu$, $\mu_n$ & — & $-\Delta_\gamma Y_n=\mu_n^2Y_n$ & Kaluza–Klein mass \\
$\vec n$, $k$ & — & $\vec n\in\mathbb Z^4$, $k=|\vec n|^2$ & torus mode number; shell \\
$r_4(k)$, $\sigma(k)$ & \texttt{rIIr4}, \texttt{rIIsig} & (7.1) & number of $\vec n$ with $|\vec n|^2=k$; divisor sum \\
$l$, $d_l$ & \texttt{rdDl[l]} & (7.2) & degree of a spherical harmonic; multiplicity \\
$N(\mu)$, $x$, $L$, $y$ & \texttt{rdNS4[x]}, \texttt{rdWeyl} & (7.3) & number of states below $\mu$; $x=\mu r$; largest $l$; $y=L+2$ \\
$Z_l$, $C_l^{(3/2)}$ & \texttt{rdZon[l]}, \texttt{GegenbauerC} & (7.4) & zonal harmonic; Gegenbauer polynomial \\
$L_{ab}$ & — & $X_a\partial_b-X_b\partial_a$ & rotations of $\mathbb R^5$ (Killing vectors of $S^4$) \\*
\end{longtable}}
\begin{align*}
r_4(k)&=8\sigma(k)-32\sigma(k/4),\quad \sigma(k)=\sum_{d|k}d\ (\sigma(k/4)=0\text{ unless }4|k);\quad T^4:\ \mu=\frac{2\pi|\vec n|}{\ell_\perp},&&(7.1) \\
S^4:\ \mu&=\frac{\sqrt{l(l+3)}}r,\quad d_l=\frac{(2l+3)(l+1)(l+2)}6=\binom{l+4}4-\binom{l+2}4,&&(7.2) \\
N(\mu)&\simeq\frac{V\mu^4}{32\pi^2};\quad S^4:\ N=\sum_{l\le L}d_l=\frac{y^2(y^2-1)}{12},\quad L(L+3)\le x^2<(L+1)(L+4),&&(7.3) \\
Z_l(x_4)&=\sqrt{\frac{d_l}V}\,\frac{C_l^{(3/2)}(\cos x_4)}{C_l^{(3/2)}(1)},\qquad C_l^{(3/2)}(1)=\frac{(l+1)(l+2)}2.&&(7.4)
\end{align*}
\Needspace{12\baselineskip}\section*{8. Spinors on \texorpdfstring{$S^4$}{S4} and dirac16complex}\addcontentsline{toc}{section}{8. Spinors on \texorpdfstring{$S^4$}{S4} and dirac16complex}
{\footnotesize
\begin{longtable}{>{\raggedright\arraybackslash}p{\dimexpr 0.1700\linewidth-2\tabcolsep\relax}>{\raggedright\arraybackslash}p{\dimexpr 0.2200\linewidth-2\tabcolsep\relax}>{\raggedright\arraybackslash}p{\dimexpr 0.3100\linewidth-2\tabcolsep\relax}>{\raggedright\arraybackslash}p{\dimexpr 0.3000\linewidth-2\tabcolsep\relax}}
\toprule
\textbf{symbol} & \textbf{Mathematica} & \textbf{definition} & \textbf{meaning} \\
\midrule\endfirsthead
\toprule
\textbf{symbol} & \textbf{Mathematica} & \textbf{definition} & \textbf{meaning} \\
\midrule\endhead
\endfoot
\bottomrule\endlastfoot
$\rho_k$ & \texttt{rhoK} (section 17.2) & $\rho_k\rho_l+\rho_l\rho_k=-2\delta_{kl}$ & the $4\times4$ matrices of the transverse spinor index \\
$\tilde\gamma_a$ & \texttt{rdGt} & $i\rho_{a+3}$ & Hermitian Dirac matrices of the extra space \\
$e^a$, $e_i{}^a$ & \texttt{rdHf} & (8.1) & orthonormal frame of $S^4$ \\
$\omega$ & \texttt{rdOm} & spin connection & connection coefficients of the frame \\
$\nabla_i$ & \texttt{rdNab} & (8.1) & spinor covariant derivative \\
$\not D$ & — & $\tilde\gamma^ae_a{}^i\nabla_i$ & Dirac operator of the extra space \\
$\varepsilon_s$, $\varepsilon_0$, $s$ & \texttt{rdKS[s]} & (8.2) & Killing spinors; a constant spinor; $s=\pm1$ \\
$\lambda$ (Dirac) & — & (8.3) & eigenvalue of $\not D$ \\
$m$ & — & mass parameter & eight-dimensional mass of dirac16complex \\*
\end{longtable}}
\begin{align*}
e^a&=(r,\ r\sin x_4,\ r\sin x_4\sin x_5,\ r\sin x_4\sin x_5\sin x_6)\,dx^{a+3},\qquad \nabla_i=\partial_i+\frac14\omega_{iab}\tilde\gamma^a\tilde\gamma^b,&&(8.1) \\
\varepsilon_s&=e^{s\frac i2x_4\tilde\gamma_1}e^{\frac12x_5\tilde\gamma_1\tilde\gamma_2}e^{\frac12x_6\tilde\gamma_2\tilde\gamma_3}e^{\frac12x_7\tilde\gamma_3\tilde\gamma_4}\varepsilon_0,\quad \nabla_i\varepsilon_s=s\frac i{2r}e_i{}^a\tilde\gamma_a\varepsilon_s,\quad \not D\varepsilon_s=s\frac{2i}r\varepsilon_s,&&(8.2) \\
\lambda&=\pm\frac{i(l+2)}r\ \text{(degeneracy }4\tbinom{l+3}3\text{ each)},\qquad |\lambda|^2\ge\frac{nR_{\rm int}}{4(n-1)}=\frac4{r^2},\qquad m_{\rm KK}=\sqrt{m^2+|\lambda|^2}.&&(8.3)
\end{align*}
\Needspace{12\baselineskip}\section*{9. The warped radion (Randall–Sundrum)}\addcontentsline{toc}{section}{9. The warped radion (Randall–Sundrum)}
{\footnotesize
\begin{longtable}{>{\raggedright\arraybackslash}p{\dimexpr 0.1700\linewidth-2\tabcolsep\relax}>{\raggedright\arraybackslash}p{\dimexpr 0.2200\linewidth-2\tabcolsep\relax}>{\raggedright\arraybackslash}p{\dimexpr 0.3100\linewidth-2\tabcolsep\relax}>{\raggedright\arraybackslash}p{\dimexpr 0.3000\linewidth-2\tabcolsep\relax}}
\toprule
\textbf{symbol} & \textbf{Mathematica} & \textbf{definition} & \textbf{meaning} \\
\midrule\endfirsthead
\toprule
\textbf{symbol} & \textbf{Mathematica} & \textbf{definition} & \textbf{meaning} \\
\midrule\endhead
\endfoot
\bottomrule\endlastfoot
$A(y)$, $k$ & \texttt{rdA[rdy]}, \texttt{rdk} & (9.1) & warp function; AdS curvature scale \\
$y$, $R$ & \texttt{rdy}, \texttt{rdRR} & $y\in[-\pi R,\pi R]$, $y\simeq-y$ & the fifth coordinate; the radius \\
$T(x)$ & \texttt{rdT[rdt]}, \texttt{rdTs} & (9.2) & the radius as a field (geometric radion) \\
$\mu$, $f$ & \texttt{rdmu}, \texttt{rdfm} & (9.2) & canonical warped radion; its normalization \\
$f_{\rm UV}$, $f_{\rm IR}$, $f$ & \texttt{rdfUV}, \texttt{rdfIR}, \texttt{rdf} & (9.1) & brane tensions \\
$\lambda_A$, $\lambda_B$, $\lambda_{\rm UV}$, $\lambda_{\rm IR}$ & — & the note’s tensions & the same tensions in the note \\
$K^\mu{}_\nu$, $S_{\mu\nu}$ & \texttt{rdKmix}, \texttt{rdKjump} & (9.1) & extrinsic curvature; brane energy–momentum \\
$M$, $M_5$, $\kappa_8$ & — & table of section 9.2 of the document & Planck scales of other normalizations \\
$\Phi(y)$, $m$, $\nu$, $\epsilon$ & \texttt{rdPhy}, \texttt{rdMg}, \texttt{rdNu}, \texttt{rdeps} & (9.3) & Goldberger–Wise scalar, its mass, exponent, small parameter \\
$v_A$, $v_B$, $C_1$, $C_2$ & \texttt{rdvA}, \texttt{rdvB}, \texttt{rdC1}, \texttt{rdC2} & (9.3) & brane values of $\Phi$; integration constants \\
$\langle\mu\rangle$, $\langle T\rangle$ & \texttt{rdmu0} & (9.4) & minimum of $V(\mu)$ \\
$c_5$, $[A']$ & \texttt{rdc5}, \texttt{rdJmp} & $3$; $2k$ & coefficient of $A''$ in $G^\mu{}_\nu$; jump of $A'$ at the UV brane \\*
\end{longtable}}
\begin{align*}
&ds^2=e^{-2A(y)}\eta_{\mu\nu}dx^\mu dx^\nu-dy^2,\quad G^\mu{}_\nu=3A''-6A'^2,\quad G^y{}_y=-6A'^2,&&(9.1) \\
&f_{\rm UV}=-f_{\rm IR}=\frac{6k}\kappa,\qquad K^\mu{}_\nu=\frac12g^{\mu\alpha}\partial_yg_{\alpha\nu}, \\
&ds^2=e^{-2kT(x)|y|}g_{\mu\nu}dx^\mu dx^\nu-T(x)^2dy^2,\quad \mu=f\,e^{-k\pi T},\quad f^2=\frac6{\kappa k},\quad M_{\rm Pl}^2=\frac{1-e^{-2\pi kR}}{\kappa k},&&(9.2) \\
&\Phi''-4k\Phi'-m^2\Phi=0,\quad \Phi=e^{2ky}(C_1e^{\nu ky}+C_2e^{-\nu ky}),\quad \nu=\sqrt{4+m^2/k^2},\quad \epsilon=\frac{m^2}{4k^2},&&(9.3) \\
&V(\mu)=4k\Big(\frac\mu f\Big)^4\Big[v_B-v_A\Big(\frac\mu f\Big)^\epsilon\Big]^2,\quad \Big(\frac{\langle\mu\rangle}f\Big)^\epsilon=\frac{v_B}{v_A},\quad m_r^2=\frac43\kappa k^2\epsilon^2v_B^2\Big(\frac{\langle\mu\rangle}f\Big)^2.&&(9.4)
\end{align*}
\Needspace{12\baselineskip}\section*{10. Numbers and constants}\addcontentsline{toc}{section}{10. Numbers and constants}
{\footnotesize
\begin{longtable}{>{\raggedright\arraybackslash}p{\dimexpr 0.1700\linewidth-2\tabcolsep\relax}>{\raggedright\arraybackslash}p{\dimexpr 0.1700\linewidth-2\tabcolsep\relax}>{\raggedright\arraybackslash}p{\dimexpr 0.3000\linewidth-2\tabcolsep\relax}>{\raggedright\arraybackslash}p{\dimexpr 0.3600\linewidth-2\tabcolsep\relax}}
\toprule
\textbf{symbol} & \textbf{Mathematica} & \textbf{value} & \textbf{source} \\
\midrule\endfirsthead
\toprule
\textbf{symbol} & \textbf{Mathematica} & \textbf{value} & \textbf{source} \\
\midrule\endhead
\endfoot
\bottomrule\endlastfoot
$M_{\rm Pl}$ & \texttt{rIIMPl} & $2.435\times10^{18}$ GeV & reduced Planck mass $(8\pi G_N)^{-1/2}$ \\
$M_8$ & \texttt{rIIM8} & $10^4$ GeV & choice of sections 22–24 \\
$m_r$ & — & 10 GeV & choice of sections 22–24 \\
$m_h$, $m_{h0}$ & \texttt{mHb} & 125.20 GeV & Particle Data Group 2025 \\
$G_F$ & \texttt{gFermi} & $1.1663788\times10^{-5}$ GeV$^{-2}$ & Particle Data Group \\
$\hbar c$ & \texttt{tbHbarCm} & $1.973269804\times10^{-16}$ GeV m & CODATA 2018 \\
38.6 \texorpdfstring{\ensuremath{\mu}}{μ}m & — & Eöt-Wash range bound for $|\alpha|=1$ & Lee et al., Phys. Rev. Lett. 124 (2020) 101101 \\
$\langle\mu\rangle$ & — & 3 TeV & the note’s warped example \\*
\end{longtable}}
\Needspace{12\baselineskip}\section*{11. Symbols that look alike}\addcontentsline{toc}{section}{11. Symbols that look alike}
\begin{itemize}
\item $\beta$: the radion field $\beta(x)$ (this document) and the constant $\beta$ of the note’s ansatz; $\beta_n$ the constant for $n$ extra dimensions. Our $\beta(x)$ equals the note’s product $\beta\phi(x)$.
\item $\alpha$: the note’s constant in $e^{2\alpha\phi}$ and the strength $\alpha$ of a Yukawa (fifth) force.
\item $\phi$, $\varphi$, $\varphi_H$, $\Phi$: the canonical radion $\phi$; the note’s dimensionless radion $\varphi$; the radial Higgs field $\varphi_H$; a bulk scalar $\Phi$.
\item $\lambda$: the Higgs quartic coupling, the Dirac eigenvalue on $S^4$, the range of a Yukawa force, and (with subscripts) the note’s brane tensions.
\item $\epsilon$, $\varepsilon$: the Goldberger–Wise parameter $\epsilon=m^2/(4k^2)$; the energy density $\varepsilon$; the Killing spinors $\varepsilon_s$.
\item $\gamma$: the reference metric $\gamma_{ij}$ ($\hat\gamma$ for the unit sphere), the Higgs–radion parameter $\gamma=v/\Lambda_r$, and the photon in $h\to\gamma\gamma$; $\tilde\gamma_a$ are Dirac matrices.
\item $\sigma$: the curvature switch of $V(\beta)$ and the divisor sum $\sigma(k)$.
\item $r$: the radius of the sphere, the distance between two masses, and the bare radion field $r(x)$; $r_m$ the radion mass eigenstate; $r_4(k)$ the lattice multiplicity.
\item $h$: the induced metric $h_{\mu\nu}$, the bare Higgs field $h(x)$ and the mass eigenstate $h_m$.
\item $R$: scalar curvature ($R$, $R_E$, $R_{\rm int}$), the Randall–Sundrum radius $R$, and $R(h)$, the curvature of the induced metric.
\item $K$: the stiffness $K_U$, the kinetic matrix $K$, the extrinsic curvature $K^\mu{}_\nu$.
\item $f$: the brane tension $f$ ($f_{\rm UV}$, $f_{\rm IR}$), the normalization $f$ of the warped radion, and $f_W$ (the tension of the world brane of 7+1).
\item $\Lambda$, $\Lambda_r$: the bulk cosmological constant; the coupling scale of the radion.
\item $M$: a point mass, the Randall–Sundrum scale $M$, $M_5$, $M_8$, $M_{\rm Pl}$, and the mass matrix $M'$.
\item $S$: the action $S$, the kinetic shift matrix $S$, the brane energy–momentum $S_{\mu\nu}$, the sphere $S^4$.
\item $x$: the coordinates and the dimensionless $x=\mu r$ of Weyl’s law; $y$: the extra coordinates, the fifth coordinate of RS1, and $y=L+2$.
\item $Z$: the normalization $Z$ and the zonal harmonics $Z_l$; $\mathbb Z$ the integers.
\item $c$: the coefficient $c=-6$ and the coefficients $c_n$; $C_1$, $C_2$, $C_h$, $C_r$, $C_l^{(3/2)}$ are different objects.
\end{itemize}
\Needspace{12\baselineskip}\section*{12. Mathematica names}\addcontentsline{toc}{section}{12. Mathematica names}
\Needspace{8\baselineskip}\subsection*{12.1 The names of section 26}\addcontentsline{toc}{subsection}{12.1 The names of section 26}
{\footnotesize
\begin{longtable}{>{\raggedright\arraybackslash}p{\dimexpr 0.4200\linewidth-2\tabcolsep\relax}>{\raggedright\arraybackslash}p{\dimexpr 0.5800\linewidth-2\tabcolsep\relax}}
\toprule
\textbf{name} & \textbf{meaning} \\
\midrule\endfirsthead
\toprule
\textbf{name} & \textbf{meaning} \\
\midrule\endhead
\endfoot
\bottomrule\endlastfoot
\texttt{rdChr}, \texttt{rdRic}, \texttt{rdScal}, \texttt{rdEin}, \texttt{rdEL}, \texttt{rdEL2} & Christoffel symbols, Ricci tensor, scalar curvature, mixed Einstein tensor of a diagonal metric; Euler–Lagrange operators in $x_0$ and in $(x_0,x_1)$ \\
\texttt{rdGam4}, \texttt{rdRicS4}, \texttt{rdR0}, \texttt{rdB0}, \texttt{rdLc} & $\hat\gamma$; the Ricci tensor of $S^4$; $r_0$; $b_0$; $L_c$ \\
\texttt{rdMini[gam, rint, v0]} & the reduction test of section 26.1 (homogeneous metrics; Euler–Lagrange expressions for $N$, $a$, $\beta$) \\
\texttt{rdLp}, \texttt{rdSa}, \texttt{rdBt}, \texttt{rdPh} & $N(x_0)$, $a(x_0)$, $\beta(x_0)$, $\varphi(x_0)$ \\
\texttt{rdKap}, \texttt{rdLam}, \texttt{rdU}, \texttt{rdU0}, \texttt{rdU1}, \texttt{rdKU}, \texttt{rdUq}, \texttt{rdMP2} & $\kappa$, $\Lambda$, $U$, $U_0$, $U_1$, $K_U$, $U(\beta)$, $M_{\rm Pl}^2$ \\
\texttt{rdKinOK}, \texttt{rdKinC}, \texttt{rdKinCs}, \texttt{rdY} & the kinetic-term test; the solved coefficient $c_n$; $(c_1,\dots,c_4)$; extra coordinates of $T^n$ \\
\texttt{rdES}, \texttt{rdET}, \texttt{rdE00}, \texttt{rdE11}, \texttt{rdE44}, \texttt{rdb} & Einstein tensors of the sphere and torus cosmologies; their components; the symbol $b$ \\
\texttt{rdSolS}, \texttt{rdRadS}, \texttt{rdEps}, \texttt{rdPa}, \texttt{rdPb} & the solution for $H_a'$, $H_b'$; the radion equation on $S^4$; $\varepsilon$, $p_a$, $p_b$ \\
\texttt{rdBe}, \texttt{rdV}, \texttt{rdRHS} & $\beta$ (number); $V(\beta)$; the right-hand side of the radion equation \\
\texttt{rdSolT}, \texttt{rdSolS2}, \texttt{rdUb}, \texttt{rdR2} & the conditions on $T^4$ and on $S^4$; $U$ at the solution; the radion equation there \\
\texttt{rdVL}, \texttt{rdLamL}, \texttt{rdHE2}, \texttt{rdM2L}, \texttt{rdDS} & the \texorpdfstring{\ensuremath{\Lambda}}{Λ}-only sphere: $V$, $\Lambda$, $H^2$, $m^2$, the field equations of $dS_4\times S^4$ \\
\texttt{rdVF}, \texttt{rdFx}, \texttt{rdSolF}, \texttt{rdStF}, \texttt{rdQ2}, \texttt{rdRhoF} & the flux sphere: $V$, $F$, the solution, the stress of $W$, $q^2$, $\rho_F$ \\
\texttt{rdGdS}, \texttt{rdSqG}, \texttt{rdPhi}, \texttt{rdMm} & metric and $\sqrt{|G|}$ of the Goldberger–Wise check; $\Phi(x_4,x_5)$; $m_8$ \\
\texttt{rdPP}, \texttt{rdMs}, \texttt{rdNn}, \texttt{rdEW} & the expansion of $-Me^{-n\beta/2}$; $M$; $n$; the Eöt-Wash numbers \\
\texttt{rdDl}, \texttt{rdZon}, \texttt{rdNS4}, \texttt{rdWeyl} & $d_l$; $C_l^{(3/2)}(\cos x_4)$; the exact count on $S^4$; Weyl deviations \\
\texttt{rdEmb}, \texttt{rdJe} & the embedding $X_a$ and its Jacobian \\
\texttt{rdGt}, \texttt{rdHf}, \texttt{rdOm}, \texttt{rdNab}, \texttt{rdKS}, \texttt{rdZeroQ}, \texttt{rdv4} … \texttt{rdv7} & $\tilde\gamma_a$; frame lengths; connection; $\nabla_i$; Killing spinors; the exact zero test (half angles) \\
\texttt{rdSpW}, \texttt{rdRS4}, \texttt{rdToday} & spinor Weyl ratio; $r_{S^4}$; the numbers today \\
\texttt{rdXs5}, \texttt{rdt}, \texttt{rdx1} … \texttt{rdx3}, \texttt{rdy}, \texttt{rdz1} … \texttt{rdz3} & coordinates of the five- and eight-dimensional warped metrics \\
\texttt{rdG5}, \texttt{rdA}, \texttt{rdc5}, \texttt{rdJmp}, \texttt{rdk}, \texttt{rdf}, \texttt{rdfUV}, \texttt{rdfIR} & RS1 Einstein tensor; $A(y)$; $3$; $2k$; $k$; a tension; $f_{\rm UV}$, $f_{\rm IR}$ \\
\texttt{rdKmix}, \texttt{rdKjump}, \texttt{rdRR} & $K^\mu{}_\nu$; its jump at the UV brane; $R$ \\
\texttt{rdR5T}, \texttt{rdT}, \texttt{rdTs}, \texttt{rdT1}, \texttt{rdT2}, \texttt{rdLb}, \texttt{rdC2}, \texttt{rdL1} & the radion reduction of RS1: curvature, $T(x_0)$, $T$, $T'$, $T''$, Lagrangian, coefficients \\
\texttt{rdG8w}, \texttt{rdA1}, \texttt{rdA2} & eight-dimensional warped Einstein tensor; $A'$, $A''$ as unknowns \\
\texttt{rdGWL}, \texttt{rdPhy}, \texttt{rdMg}, \texttt{rdGWsol}, \texttt{rdC1}, \texttt{rdNu} & Goldberger–Wise Lagrangian; $\Phi(y)$; $m$; the solution; $C_1$; $\nu$ \\
\texttt{rdVgw}, \texttt{rdmu}, \texttt{rdfm}, \texttt{rdvA}, \texttt{rdvB}, \texttt{rdeps}, \texttt{rdmu0}, \texttt{rdV2gw}, \texttt{rdMgw}, \texttt{rdMgwNote} & $V(\mu)$, $\mu$, $f$, $v_A$, $v_B$, $\epsilon$, $\langle\mu\rangle$, $V''(\langle\mu\rangle)$, the mass coefficient with $f^2=6/(\kappa k)$ and with $12/(\kappa k)$ \\
\texttt{rdc4}, \texttt{rdOmg}, \texttt{rdOmE}, \texttt{rdO1}, \texttt{rdO2}, \texttt{rdEp} & $c$; $\Omega(x_0,x_1)$; $\Omega=1+o_1\epsilon r+o_2\epsilon^2r^2$; $1/\Lambda_r$; $o_2$; bookkeeping $\epsilon$ \\
\texttt{rdRr}, \texttt{rdHh}, \texttt{rdPhE}, \texttt{rdVv}, \texttt{rdRh}, \texttt{rdLbr}, \texttt{rdLamH}, \texttt{rdXi} & $r$, $h$, $\varphi_H$, $v$, $R(h)$, the brane Lagrangian, $\lambda$, $\xi$ \\
\texttt{rdMr0}, \texttt{rdMh0}, \texttt{rdHL1}, \texttt{rdHL2}, \texttt{rdGam}, \texttt{rdLq} & $m_{r0}$, $m_{h0}$; first- and second-order Lagrangians; $\gamma$; the quadratic Lagrangian \\
\texttt{rdKhr}, \texttt{rdKrr}, \texttt{rdKm}, \texttt{rdKnote}, \texttt{rdSm}, \texttt{rdZ}, \texttt{rdZ2} & kinetic coefficients and matrix; the note’s literal matrix; $S$; $Z$; $Z^2$ \\
\texttt{rdMp}, \texttt{rdQ}, \texttt{rdth}, \texttt{rdHRnum}, \texttt{rdHR71} & $M'$; $Q(\theta)$; $\theta$; the numerical mixing (60 digits); its values for 7+1 \\*
\end{longtable}}
\Needspace{8\baselineskip}\subsection*{12.2 Names of the earlier sections that section 26 uses}\addcontentsline{toc}{subsection}{12.2 Names of the earlier sections that section 26 uses}
{\footnotesize
\begin{longtable}{>{\raggedright\arraybackslash}p{\dimexpr 0.3800\linewidth-2\tabcolsep\relax}>{\raggedright\arraybackslash}p{\dimexpr 0.1000\linewidth-2\tabcolsep\relax}>{\raggedright\arraybackslash}p{\dimexpr 0.5200\linewidth-2\tabcolsep\relax}}
\toprule
\textbf{name} & \textbf{section} & \textbf{meaning} \\
\midrule\endfirsthead
\toprule
\textbf{name} & \textbf{section} & \textbf{meaning} \\
\midrule\endhead
\endfoot
\bottomrule\endlastfoot
\texttt{X}, \texttt{ID4} & 1 & the coordinates $\{x_0,\dots,x_7\}$; the $4\times4$ unit matrix \\
\texttt{curvature[]}, \texttt{gCos}, \texttt{cvCos}, \texttt{a3[x0]}, \texttt{b4[x0]} & 11 & the curvature routine; the cosmological metric and its curvature \\
\texttt{toH}, \texttt{Ha}, \texttt{Hb}, \texttt{dHa}, \texttt{dHb} & 11 & rewriting $a'$, $a''$, $b'$, $b''$ in terms of $H_a$, $H_b$ and their derivatives \\
\texttt{feStress}, \texttt{bS} & 20 & the stress tensor $(\varepsilon,p_a,p_b)$ of a comoving energy $W(a,b)$; the symbol $b$ \\
\texttt{rhoK} & 17.2 & $\rho_k$, the transverse $4\times4$ matrices \\
\texttt{vEW}, \texttt{gFermi}, \texttt{mHb} & 17.3 & $v$, $G_F$, $m_h$ \\
\texttt{rIIMPl}, \texttt{rIIM8}, \texttt{rIIVp}, \texttt{rIIr4} & 22.2 & $M_{\rm Pl}$, $M_8$, $V_\perp$, $r_4(k)$ \\
\texttt{tbLp}, \texttt{tbHbarCm}, \texttt{tbNmu} & 25 & $\ell_\perp$, $\hbar c$, $N(\mu)$ on $T^4$ \\
\texttt{check} & 0 & records a named test in the association \texttt{checks} \\*
\end{longtable}}
\end{document}
